% For now, this is a seperate document.
% However, if more components of Atmlab grow documentation in the future, this
% may become part of a larger body of Atmlab documentation.
%\documentclass[a4paper,10pt]{memoir}
\documentclass[a4paper,10pt]{article}
\usepackage{standalone}
\usepackage[T1]{fontenc}
\usepackage{lmodern}
\usepackage{extsizes} % redundant with memoir
\usepackage{natbib}
\setcitestyle{authoryear,round,semicolon}
\usepackage{paralist}
\usepackage[colorlinks=true,allcolors=black]{hyperref}
\usepackage[style=index, acronym]{glossaries}
\usepackage{microtype}
\usepackage{amsmath}
\usepackage{amssymb}
\usepackage{textcomp}
\usepackage{siunitx}
\sisetup{seperr,repeatunits=false}
\usepackage{booktabs}
\usepackage{tikz}
\usepackage{pgfplots}
\pgfplotsset{compat=1.5}
%\usetikzlibrary{arrows,positioning,backgrounds,fit,shapes,calc}
%\usepgfplotslibrary{units}
\hbadness=10000 % don't bother be about underfull hbox
\vbadness=10000 % idem dito

\usepackage[utf8]{inputenc}
\usepackage[english]{babel}
\usepackage[english]{isodate}
\usepackage{url}
\usepackage{pgf-umlcd}
\makeglossaries
% For use with \usepackage[acronym]{glossaries}.
%
% \input this file after you have defined your own acronyms (but why not
% add your own acronyms here?). It will define any acronyms that you have
% yet defined yourself (acronyms will not be over-written). The acronym
% key will be equal to the actual acronym, e.g. \gls{ARTS}.
%
% You don't need to include an actual glossary to use those acronyms; the
% package takes care of the first usage \gls{ABC} being expanded. See
% 'texdoc glossaries' for more info.
%
% If you do include a glossary (table of acronyms), only those acronyms
% that are actually used will be included.
%
% $Id: acronyms.tex 7678 2012-04-04 16:30:11Z gerrit $

\newcommand{\makeacroif}[2]{%
\ifglsentryexists{#1}{}{%
 \newacronym{#1}{#1}{#2}}}

\newcommand{\makepluralacroif}[3]{%
\ifglsentryexists{#1}{}{%
 \newacronym[plural=#2]{#1}{#1}{#3}}}

\makeacroif{IPCC}{Intergovernmental Panel on Climate Change}
\makeacroif{AR4}{4th Assessment Report}
\makeacroif{NOAA}{National Oceanic and Atmospheric Administration}
\makeacroif{TRMM}{Tropical Rainfall Measuring Mission}
\makeacroif{CALIPSO}{Cloud-Aerosol Lidar and Infrared Pathfinder Satellite Observation}
\makeacroif{PARASOL}{Polarization and Anisotropy of Reflectances
                              for Atmospheric Sciences Coupled with
                              Observations from a Lidar}
\makeacroif{OCO}{Orbital Carbon Observatory}
\makeacroif{ESMR}{Electrically Scanning Microwave Radiometer}
\makeacroif{DMSP}{Defense Meteorological Satellite Program}
\makeacroif{SSM/T}{Special Sensor Microwave/Temperature}
\makeacroif{SEVIRI}{Spinning Enhanced Visible Infra-Red Imager}
\makeacroif{VHRR}{Very High Resolution Radiometer}
\makeacroif{UV}{Ultra-Violet}
\makeacroif{IR}{Infra-Red}
\makeacroif{MW}{microwave}
\makeacroif{FRAC}{Full Resolution Area Coverage}
\makeacroif{GAC}{Global Area Coverage}
\makeacroif{IWSSM}{International Workshop on Space-based Snowfall Measurement}
\makeacroif{DDA}{Discrete Dipole Approximation}
\makeacroif{PND}{particle number density}
\makeacroif{FDTD}{Finite Difference Time Domain}
\makeacroif{DOIT}{Discrete Ordinate ITerative}
\makeacroif{PR}{Precipitation Radar}
%
\makeacroif{HIRS}{High-resolution Infrared Radiation Sounder}
\makeacroif{AVHRR}{Advanced Very High Resolution Radiometer}
\makeacroif{MC}{Monte Carlo}
\makeacroif{NOAA}{National Oceanic and Atmospheric Administration}
\makeacroif{ARTS}{Atmospheric Radiative Transfer Simulator}
\makepluralacroif{mpbs}{MPBS}{monochromatic pencil-beam simulation}
\makeacroif{RMSE}{root mean square error}
\makeacroif{GCM}{general circulation model}
\makeacroif{CPR}{Cloud Profiling Radar}
\makeacroif{IWP}{Ice Water Path}
\makeacroif{AMSU}{Advanced Microwave Sounding Unit}
\makeacroif{MHS}{Microwave Humidity Sounder}
\makeacroif{ANN}{Artificial Neural Network}
\makeacroif{radar}{RAdio Detection and Ranging}
\makeacroif{IR}{infrared}
\makeacroif{MW}{microwave}
\makeacroif{MODIS}{MODerate resolution Imaging Spectroradiometer}
\makeacroif{FRAC}{Full-Resolution Area Coverage}
\makeacroif{GAC}{Global Area Coverage}
\makeacroif{TIROS}{Television InfraRed Operational Sounder}
\makeacroif{SRF}{spectral response function}
\makeacroif{DOIT}{Discrete Ordinate ITerative}
\makeacroif{DOM}{Discrete Ordinate Method}
\makeacroif{ECMWF}{European Centre for Medium-range Weather Forecasting}
\makeacroif{FASCOD}{Fast Atmospheric Signature CODe}
\makeacroif{MSU}{Microwave Sounding Unit}
\makeacroif{SSM/I}{Special Sensor Microwave/Imager}
\makeacroif{TIROS}{Television Infra-Red Operational Sounder}
\makeacroif{GCM}{General Circulation Model}
\makeacroif{ERA}{\protect\gls{ECMWF} Re-Analysis}
\makeacroif{NESDIS}{National Environment Satellite, Data and Information Service}
\makeacroif{MSPPS}{Microwave Surface and Precipitation Products System}
\makeacroif{UTH}{Upper Tropospheric Humidity}
\makeacroif{PWV}{Precipitable Water Vapour}
\makeacroif{TOA}{Top Of Atmosphere}
\makeacroif{IWC}{Ice Water Content}
\makeacroif{NPP}{National Partnership Program}
\makeacroif{EOS}{Earth Observing System}
\makeacroif{ATMS}{Advanced Technology Microwave Sounder}
\makeacroif{CrIS}{Cross-track Infrared Sounder}
\makeacroif{AIRS}{Atmospheric Infrared Sounder}
\makeacroif{FOV}{Field Of View}
\makeacroif{FOR}{Field Of Regard}
\makeacroif{VIIRS}{Visible Infrared Imager Radiometer Suite}
\makeacroif{FY}{Feng-Yun}
\makeacroif{MWHS}{Micro-Wave Humidity Sounder}
\makeacroif{MWS}{Micro-Wave Sounder}
\makeacroif{SG}{Second Generation}
\makeacroif{STAR}{Center for Satellite Applications and Research}
\makeacroif{MIRS}{Microwave Integrated Retrieval System}
\makeacroif{TIROS}{Television InfraRed Operational Satellite}
\makeacroif{TOVS}{\protect\gls{TIROS} Operational Vertical Sounder}
\makeacroif{ATOVS}{Advanced \protect\gls{TIROS} Operational Vertical Sounder}
\makeacroif{AAPP}{Advanced \protect\gls{ATOVS} Processing Package}
\makeacroif{POES}{Polar Orbiting Environmental Satellites}
\makeacroif{NASA}{National Aeronautics and Space Administration}
\makeacroif{HIRLAM}{High Resolution Local Area Modelling}
\makeacroif{NWP}{Numerical Weather Prediction}
\makeacroif{HARMONIE}{\protect\gls{HIRLAM} Aladin Regional/Mesoscale Operational \protect\gls{NWP} In Europe}
\makeacroif{RTSP}{Radiative Transfer and Surface Properties}
\makeacroif{WG}{Working Group}
\makeacroif{ITWG}{International \protect\gls{TOVS} Working Group}
\makeacroif{RDR}{Raw Data Record}
\makeacroif{SDR}{Sensor Data Record}
\makeacroif{CDR}{Climate Data Record}
\makeacroif{VTPR}{Vertical Temperature Profile Radiometer}
\makeacroif{DWD}{Deutsche Wetter-Dienst}
\makeacroif{JPL}{Jet Propulsion Library}
\makeacroif{SSU}{Stratospheric Sounding Unit}
\makeacroif{C3VP}{Canadian CloudSat/CALIPSO Validation Project}
\makeacroif{SNO}{Simultaneous Nadir Observation}
\makeacroif{NRL}{Naval Research Laboratory}
\makeacroif{DARDAR}{raDAR/liDAR}
\makeacroif{CALIOP}{Cloud-Aerosol Lidar with Orthogonal Polarization}
\makeacroif{VLIDORT}{Vectorized Linearized Discrete Ordinate Radiative Transfer}
\makeacroif{DISORT}{Discrite Ordinate Radiative Transfer}
\makeacroif{LMD}{Laboratoire de M\'{e}t\'{e}orologie Dynamique}
\makeacroif{CLEX}{Cloud Layer EXperiment}
\makeacroif{CSU}{Colorado State University}
\makeacroif{IIR}{Imaging Infrared Radiometer}
\makeacroif{EUMETSAT}{European Organisation for the Exploitation of Meteorological Satellites}
\makeacroif{CM}{Climate Monitoring}
\makeacroif{SAF}{Satellite Application Facility}
\makeacroif{TWP}{Total Water Path}
\makeacroif{LERMA}{Laboratoire d'Etude du Rayonnement et de la Mati\`ere en Astrophysique}
\makeacroif{ITSC}{International \protect\gls{TOVS} Study Conference}
\makeacroif{PSD}{Particle Size Distribution}
\makeacroif{FASTEM}{Fast Emissivity (FIXME CHECK ACRONYM!)}
\makeacroif{TELSEM}{A Tool to Estimate Land Surface Emissivities at Microwave}
\makeacroif{RTTOV}{Radiative Transfer for \protect\gls{TOVS}}
\makeacroif{ASG}{Atmospheric Scene Generator}
\makeacroif{NICAM}{Nonhydrostatic ICosahedral Atmospheric Model}


\title{Collocation Toolkit in Atmlab}

\date{\isodate\today}

\author{Gerrit Holl}

\begin{document}
\pagenumbering{Roman}
\maketitle
%\newpage
\tableofcontents
%\newpage
\listoffigures
%\newpage
\listoftables
%\newpage
\pagenumbering{arabic}

\section{Introduction}

A frequent task in analysis of geographical data is the collocation in space
and time between data from different datasets.
For example, one might wish to collocate geolocated remote sensing measurements from
spaceborne sensors with model output, or collocate measurements from different
sensors with each other.
Usually, someone develops code in his or her favourite programming language
to tackle the collocation problem at hand.
Depending on context, the problem can vary considerably in nature and in
complexity.
Nevertheless, many sub-problems occur in a wide variety of situations.
Valuable time may be lost by people re-inventing the wheel over and over
again.

The Collocations Toolkit is a set of classes and routines to easily
handle collocations between an arbitrary pair of (satellite) datasets.
It has been designed for user-friendly processing of long time series of data in a
wide variety of situations.

The toolkit is currently implemented in Matlab and is fully integrated with
Atmlab, a package for atmospheric data analysis.
It is available at a free license and is completely free of charge.

Practical information on how to download the code can be found at
\url{http://www.sat.ltu.se/projects/collocations/toolkit.php}.

This is a technical document, and does not contain detailed information on
algorithms used in the process of collocating.
For information on the algorithms, please refer to the descriptions in the
refereed articles \citet{holl10:_collocating_amt} and
\citet{john12:_understanding_jgr}, and in my licentiate thesis
\citep{holl11:_lic}.

\section{Using the Toolkit}

The diagram below illustrates the most important classes and methods.

\begin{tikzpicture}[every node/.style={font=\scriptsize}]
 \begin{class}{SatDataset}{0, 0}
  \attribute{+ name}
  \attribute{+ basedir}
  \attribute{+ subdir}
  \attribute{+ re}
  \attribute{+ filename}
  \attribute{+ reader : function\_handle}
  \attribute{+ granule\_duration : double}
  \attribute{+ satname}
  \attribute{+ dataset}
  \attribute{+ firstline\_filename}
  \attribute{\# collocated\_dataset}
  \operation{+ find\_datadir\_by\_date}
  \operation{+ find\_info\_from\_granule}
  \operation{+ find\_granules\_by\_date}
  \operation{+ find\_granule\_by\_datetime}
  \operation{+ read\_granule}
  \operation{+ granule\_first\_line}
  \operation{+ find\_granule\_first\_line}
 \end{class}

 \begin{class}{HomemadeDataset}{1, -7}
  \inherit{SatDataset}
  \attribute{+ cols}
  \operation{+ store}
  \operation{\# read\_single\_day}
 \end{class}

 \begin{class}{CollocatedDataset}{-5, -10}
 \inherit{HomemadeDataset}
 \attribute{+ primary}
 \attribute{+ secondary}
 \attribute{\# associated}
 \attribute{\# members}
 \operation{+ overlap\_granule}
 \operation{+ collocate\_granule}
 \operation{+ collocate\_date}
 \operation{+ collocate\_and\_store\_date}
 \operation{+ collocate\_and\_store\_date\_range}
 \operation{+ process}
 \operation{+ read}
 \end{class}

 \begin{abstractclass}{AssociatedDataset}{3, -10}
 \inherit{HomemadeDataset}
 \attribute{+ dependencies}
 \operation{+ members2cols}
 \operation{+ process\_delayed}
 \operation{+ merge}
 \operation{+ get\_mergefields}
 \operation{+ concatenate}
 \operation[0]{+ primary\_arguments}
 \operation[0]{+ secondary\_arguments}
 \operation[0]{+ needs\_primary\_data}
 \operation[0]{+ needs\_secondary\_data}
 \operation[0]{+ process\_granule}
 \end{abstractclass}

 \begin{class}{FieldCopier}{-2, -15}
 \implement{AssociatedDataset}
 \end{class}

 \begin{class}{MeanFieldCopier}{5, -15}
 \implement{AssociatedDataset}
 \end{class}

 \unidirectionalAssociation{MeanFieldCopier}{averages}{}{FieldCopier}
 \association{CollocatedDataset}{consists of two}{}{SatDataset}{can have many}{}


\end{tikzpicture}

\section{Developing the Toolkit}

\printglossaries

\bibliographystyle{copernicus}
\bibliography{bib_all}

\end{document}
